\section{The Relative Zeta Value}\label{an:section}

We bound the relative zeta value $L_F(1)$ needed in the geometric
construction. Logarithms of Dedekind zeta functions are normalized by the
degree: for a number field $A$ we work with $[A:\Q]^{-1}\log\zeta_A(s)$.

\begin{definition}\label{an:family}
Let $\mathcal K$ be the set of all finite Galois extensions $K/B$ contained
in $B_G$, of degree a power of two at least $2^{49}$, such that
$G_B\to\overline G_B$ factors through $\Gal(K/B)$ (notation of
Definition~\ref{tw:GB}). These are the fields $K$ of
Theorem~\ref{tw:field-family}.
\end{definition}

For $K\in\mathcal K$ we use the notation of
Theorem~\ref{tw:field-family}: $\iota_1$ is a complex conjugation at a place of $K$ above
$v_1$, $F=K^{\langle\iota_1\rangle}$, $d=[F:\Q]$, so that $[K:\Q]=2d$, $(b,c)$
is the signature of $F$, $\theta=c/d$, and $\lambda=2^{9/4}\sqrt{3615}$ is the
root discriminant of both $K$ and $F$. Put
\[
 \ell=\log\lambda=\tfrac94\log2+\tfrac12\log3615=5.65600472355\ldots
\]
Different choices of $\iota_1$ are conjugate in $\Gal(K/B)$ and give
conjugate fields $F$; hence $b$, $c$, $\theta$ and the function
$L_F(s)=\zeta_K(s)/\zeta_F(s)$ depend only on $K$. The extension $K/F$ is
quadratic and unramified at every finite place, $K$ is totally imaginary,
and $b\ge1$. Let $\gamma$ be Euler's constant. For a
prime $\mathfrak p$ of $B$, let $e_K(\mathfrak p)$ and $f_K(\mathfrak p)$ be
its ramification index and residue degree in the Galois extension $K/B$, so
that $(e_K(\mathfrak p),f_K(\mathfrak p))$ is the type of $\mathfrak p$ in
$K/B$ (Definition~\ref{tw:type}).


\begin{theorem}\label{an:ceiling}
There is an $m_0$ such that every $K\in\mathcal K$ with $[K:B]\ge m_0$
satisfies
\[
 \frac1d\log L_F(1)<C_{\rm eff}:=\Ceff.
\]
\end{theorem}

The proof of Theorem~\ref{an:ceiling} is completed in
Section~\ref{ce:section}. This section proves the general criterion,
Proposition~\ref{an:signed}, and encloses $\zeta_E$ at real points
(Proposition~\ref{an:genus-value}); Sections~\ref{dh:section}
and~\ref{lv:section} supply the other inputs of the criterion.

The proof assumes neither the generalized Riemann hypothesis nor a
Brauer--Siegel type asymptotic formula, and it does not give an effective
value of $m_0$. It has four steps. Monotonicity of the completion of $L_F$
reduces the value at $s=1$ to $\zeta_K$ at real points $\sigma>1$
(Subsection~\ref{an:reduction}). Along a sequence of fields of growing
degree, the normalized numbers of primes of each prime power norm $q$
converge to limiting densities $\beta_q$ (Lemma~\ref{an:limit}), and
$d^{-1}\log L_F(1)$ is bounded in the limit by $Z(1)=\sum_q\beta_qg_1(q)$,
where $g_s(q)=-\log(1-q^{-s})$; the Tsfasman--Vl\u{a}du\c{t} inequality
bounds a weighted sum of the $\beta_q$ (Subsection~\ref{an:budget-section}).
The \emph{selected primes} of $B$, those above $2$, $3$, $5$ and the capped
primes, have the same type in every field of the tower, so their part of
$Z(1)$ is known exactly (Subsection~\ref{an:selected}). For each other prime
of $B$ we use a lower bound for its residue degree, a \emph{floor}, and a
kernel inequality with several abscissae bounds the rest of $Z(1)$ by a
combination, with coefficients of both signs, of \emph{floor sums}: the sums
of the terms of these primes at residue degree equal to the floor, evaluated
at real points $\sigma_j>1$ (Proposition~\ref{an:signed}).
Sections~\ref{dh:section}--\ref{ce:section} supply the floors and the
values: the floors come from a field $\EW$ contained in every field of the
tower and from a census, and the values come from the factorization of
$\zeta_{\EW}$ into $L$-functions. Throughout, $\sigma$ denotes a real number
greater than $1$ and $\varepsilon=\sigma-1$; in this section $\varepsilon$ is
this positive number, not the unit $\alpha_1$ of \eqref{tw:alpha}.

\subsection{\texorpdfstring{Reduction to $s>1$}{Reduction to s > 1}}
\label{an:reduction}
This subsection shows that $L_F(1)>0$, by the class-number formula, and that
the completion $\mathcal L_F$ of $L_F$, defined in \eqref{an:completion}
below, is positive and nondecreasing on
$[1,\infty)$ (Lemma~\ref{an:monotonicity}). Together with a comparison of
Euler products they bound $L_F(1)$ by $\zeta_K$ at a real point
$1+\varepsilon$ (Lemma~\ref{an:shift}), which Lemma~\ref{an:limit} uses.

For a number field $A$, write $h_A$, $\operatorname{Reg}_A$, $w_A$ and
$\Delta_A$ for its class number, regulator, number of roots of unity and
discriminant. The residue of $\zeta_A$ at $s=1$ is
$2^{r_1(A)}(2\pi)^{r_2(A)}h_A\operatorname{Reg}_A/(w_A|\Delta_A|^{1/2})$.
Since $K/F$ is unramified at the finite places, $|\Delta_K|=|\Delta_F|^2$.
Hence
\begin{equation}\label{an:class-formula}
 L_F(1)=\lim_{s\to1}\frac{\zeta_K(s)}{\zeta_F(s)}
 =2^c\pi^{d-c}\lambda^{-d/2}\,
 \frac{h_K\operatorname{Reg}_Kw_F}{h_F\operatorname{Reg}_Fw_K}>0.
\end{equation}

Let $s>1$ be real. At a prime of $F$ of norm $q$, the Euler factor of $L_F$
is $(1-q^{-s})^{-1}$ if the prime splits in $K$ and $(1+q^{-s})^{-1}$ if it
is inert. Its square is at most the Euler factor of $\zeta_K$, which is
$(1-q^{-s})^{-2}$, respectively $(1-q^{-2s})^{-1}$. Consequently
\begin{equation}\label{an:relative-euler}
 \frac{\log L_F(s)}d\le\frac{\log\zeta_K(s)}{2d}
 =\frac{\log\zeta_K(s)}{[K:\Q]}.
\end{equation}

The function $L_F$ is the Hecke $L$-function of the quadratic character of
$K/F$. This character is unramified at every finite place and nontrivial at
every real place of $F$, since $K$ is totally imaginary. With $\Gamma_\R$,
$\Gamma_\C$ and $\Lambda_A$ as in Section~\ref{gf:section}, put
\begin{equation}\label{an:completion}
 \mathcal L_F(s)=\frac{\Lambda_K(s)}{\Lambda_F(s)}
 =\lambda^{ds/2}\,\Gamma_\R(s+1)^b\,\Gamma_\C(s)^c\,L_F(s).
\end{equation}
By Hecke's theorem \cite{Hecke1920,Tate1967}, $\mathcal L_F$ is entire of
order at most one, and it satisfies $\mathcal L_F(s)=\mathcal L_F(1-s)$
because both completed Dedekind zeta functions do.

\begin{lemma}\label{an:monotonicity}
The function $\mathcal L_F(s)$ is positive and nondecreasing on $[1,\infty)$.
\end{lemma}

\begin{proof}
The Euler product shows that $\mathcal L_F$ has no zeros with
$\operatorname{Re}s>1$, and the functional equation excludes
$\operatorname{Re}s<0$. So every zero $\rho$ satisfies $0\le\operatorname{Re}
\rho\le1$. There is no zero on $[1,\infty)$, by the Euler product and by
\eqref{an:class-formula}. As $\mathcal L_F(s)>0$ for large real $s$,
it is positive on $[1,\infty)$.

Put $\Xi(z)=\mathcal L_F(\frac12+z)$. It is entire of order at most one,
even, and real on the real axis. Its zeros $z=\rho-\frac12$ satisfy
$|\operatorname{Re}z|\le\frac12$, and $\sum_{z\ne0}|z|^{-2}<\infty$. Group
the nonzero zeros into pairs $\{z_0,-z_0\}$, and let $m\ge0$ be the order of
$\Xi$ at $0$. By Hadamard's theorem for functions of order at most one,
\[
 \frac{\Xi'}{\Xi}(z)-\frac mz-\sum_{\{z_0,-z_0\}}\frac{2z}{z^2-z_0^2}
\]
is constant, where the series converges absolutely and locally uniformly
away from the zeros. Since $\Xi$ is even, $\Xi'/\Xi$ is odd, so the constant
is zero. For real $x\ge\frac12$ the logarithmic derivative is real, so it
equals the sum of the real parts. If $z_0=a+iy$, then
\[
 \operatorname{Re}\frac{2x}{x^2-z_0^2}
 =\frac{2x(x^2-a^2+y^2)}{(x^2-a^2+y^2)^2+4a^2y^2}\ge0,
\]
because $|a|\le\frac12\le x$. Since also $m/x\ge0$, the logarithmic
derivative is nonnegative, and $\log\Xi$ is nondecreasing on
$[\frac12,\infty)$.
\end{proof}

For $0<\varepsilon\le1$ and $0\le\theta\le\frac12$ put
\begin{equation}\label{an:gamma}
 \Upsilon(\varepsilon,\theta)
 =\frac\varepsilon2(\ell-\log\pi)-\varepsilon\theta\log2
 +(1-2\theta)\log\Gamma\Bigl(1+\frac\varepsilon2\Bigr)
 +\theta\log\Gamma(1+\varepsilon).
\end{equation}

\begin{lemma}\label{an:shift}
Let $K\in\mathcal K$ and $0<\varepsilon\le1$. Then
\begin{equation}\label{an:finite-gamma}
 \frac{\log L_F(1)}d
 \le\frac{\log\zeta_K(1+\varepsilon)}{[K:\Q]}+\Upsilon(\varepsilon,\theta).
\end{equation}
Moreover $\Upsilon(\varepsilon,\theta)\to0$ as $\varepsilon\downarrow0$,
uniformly for $0\le\theta\le\frac12$.
\end{lemma}

\begin{proof}
By Lemma~\ref{an:monotonicity},
$\mathcal L_F(1)\le\mathcal L_F(1+\varepsilon)$. In \eqref{an:completion}
the gamma factors change by
\[
 \frac{\Gamma_\R(2+\varepsilon)}{\Gamma_\R(2)}
 =\pi^{-\varepsilon/2}\,\Gamma\Bigl(1+\frac\varepsilon2\Bigr),\qquad
 \frac{\Gamma_\C(1+\varepsilon)}{\Gamma_\C(1)}
 =(2\pi)^{-\varepsilon}\,\Gamma(1+\varepsilon).
\]
Taking logarithms, dividing by $d$, and using $b/d=1-2\theta$, $c/d=\theta$
and \eqref{an:relative-euler}, we obtain \eqref{an:finite-gamma}. Since
$\Upsilon$ is affine in $\theta$, $\Upsilon(\varepsilon,\theta)\to0$ as
$\varepsilon\downarrow0$, uniformly for $0\le\theta\le\frac12$.
\end{proof}

The term $\Upsilon$ does not appear in the final estimate. With
$\varepsilon=1/300$, \eqref{an:finite-gamma} would add
$\Upsilon(1/300,\theta)\approx0.0054$ for $\theta$ near $\frac12$. Instead,
the proof of Lemma~\ref{an:limit} applies \eqref{an:finite-gamma} with $s-1$
in place of $\varepsilon$ for each fixed $s\in(1,2]$, passes to the limit
along a sequence of fields, and then lets $s\downarrow1$, so that $\Upsilon$
vanishes in the limit; the change from the real points $\sigma>1$ back to
$1$ is then paid for through the Tsfasman--Vl\u{a}du\c{t} inequality
(Propositions~\ref{an:prime-budget} and~\ref{an:signed}). This limit is
also the reason why $m_0$ is not effective.


\subsection{\texorpdfstring{The Tsfasman--Vl\u{a}du\c{t} Inequality over $B$}{The Tsfasman-Vladut Inequality over B}}
\label{an:budget-section}
This subsection passes to limits along sequences of fields of $\mathcal K$
(Lemma~\ref{an:limit}). The limit densities $\beta_q$ and the function $Z(s)$
defined there are the quantities that Propositions~\ref{an:prime-budget}
and~\ref{an:signed} bound, and the Tsfasman--Vl\u{a}du\c{t} inequality bounds
$\sum_q\beta_qw(q)$. Lemma~\ref{an:local-contribution} first writes the
normalized logarithm of a Dedekind zeta function as a sum over the primes
of $B$.

For real $q>1$ and $s\ge1$ put $g_s(q)=-\log(1-q^{-s})$.

\begin{lemma}\label{an:local-contribution}
Let $M$ be a finite Galois extension of $B$ and $s>1$. For a prime
$\mathfrak p$ of $B$ let $e_M(\mathfrak p)$ and $f_M(\mathfrak p)$ be its
ramification index and residue degree in $M/B$. Then
\begin{align}
 \frac{\log\zeta_M(s)}{[M:\Q]}
 &=\sum_{\mathfrak p}\frac{g_s(\mathrm N\mathfrak p^{f_M(\mathfrak p)})}
 {2e_M(\mathfrak p)f_M(\mathfrak p)},\label{an:local-sum}\\
 -\frac1{[M:\Q]}\frac{\zeta_M'}{\zeta_M}(2)
 &=\frac1{[M:\Q]}\sum_{\mathfrak P}\frac{\log\mathrm N\mathfrak P}
 {\mathrm N\mathfrak P^2-1}
 =\sum_{\mathfrak p}\frac{\log\mathrm N\mathfrak p}
 {2e_M(\mathfrak p)(\mathrm N\mathfrak p^{2f_M(\mathfrak p)}-1)},
 \label{an:local-debit}
\end{align}
where $\mathfrak P$ runs over the primes of $M$. Each summand on the right
does not increase when $e_M(\mathfrak p)$ or $f_M(\mathfrak p)$ is replaced
by a larger real number. If $M\subseteq M'$ are both Galois over $B$, then
$e_M(\mathfrak p)\mid e_{M'}(\mathfrak p)$ and
$f_M(\mathfrak p)\mid f_{M'}(\mathfrak p)$.
\end{lemma}

\begin{proof}
Above $\mathfrak p$ there are $[M:B]/(ef)$ primes of $M$, each of norm
$\mathrm N\mathfrak p^f$, and $[M:\Q]=2[M:B]$. This gives both formulas.
The first equality in \eqref{an:local-debit} follows from
$-\zeta_M'/\zeta_M(s)=\sum_{\mathfrak P}\log\mathrm N\mathfrak P/(\mathrm
N\mathfrak P^s-1)$. For monotonicity in $f$, write
$g_s(x^f)/f=\sum_{j\ge1}x^{-fjs}/(fj)$ for $x>1$; each term decreases in
$f$. The divisibilities hold because ramification indices and residue
degrees are multiplicative in towers.
\end{proof}

Put
\[
 \kappa_\infty=\frac{\ell-\gamma-\log(4\pi)}4,\qquad
 w(q)=2\log q\sum_{m\ge1}\frac1{q^m+1}.
\]
Here $2\kappa_\infty$ is the right side, and $w(q)$ the weight of the prime
power $q$, in the Tsfasman--Vl\u{a}du\c{t} inequality \eqref{an:tv-budget}
below.

\begin{lemma}\label{an:limit}
Let $C'\in\R$, and let $K^{(j)}\in\mathcal K$, $j\ge1$, be fields of strictly
increasing degree. Let $F^{(j)}$ and $d_j=[F^{(j)}:\Q]$ be the field $F$ and
the number $d$ attached to $K^{(j)}$ above, so that $[K^{(j)}:\Q]=2d_j$, and
suppose that $d_j^{-1}\log L_{F^{(j)}}(1)\ge C'$ for every $j$. For a prime
power $q$ let $N_q(K)$ be the number of
primes of $K$ of norm $q$. After passing to a subsequence, the limit
$\beta_q=\lim_jN_q(K^{(j)})/(2d_j)$ exists for every prime power $q$, and
$Z(s)=\sum_q\beta_qg_s(q)$, for real $s\ge1$, satisfies:
\begin{enumerate}
\item[(a)] $Z(s)=\lim_j[K^{(j)}:\Q]^{-1}\log\zeta_{K^{(j)}}(s)$ for every
 real $s>1$;
\item[(b)] $\sum_q\beta_qw(q)\le2\kappa_\infty$;
\item[(c)] $C'\le Z(1)<\infty$.
\end{enumerate}
\end{lemma}

\begin{proof}
If $q=p^f$, then $N_q(K)\le[K:\Q]/f$, since the local degrees above $p$ add
up to $[K:\Q]$. So $N_q(K^{(j)})/(2d_j)\le1$, and a diagonal argument gives
a subsequence along which $\beta_q$ exists for every $q$. The fields
$K^{(j)}$ have strictly increasing degrees, so they are pairwise
non-isomorphic.

(a) For fixed real $s>1$,
$(2d_j)^{-1}\log\zeta_{K^{(j)}}(s)=\sum_q(N_q(K^{(j)})/2d_j)g_s(q)$ tends to
$Z(s)$ by dominated convergence: for $q=p^f$ the terms are at most
$g_s(q)/f\le2q^{-s}$, and $\sum_qq^{-s}<\infty$.

(b) Tsfasman and Vl\u{a}du\c{t} call a sequence $(A_i)$ of pairwise
non-isomorphic number fields asymptotically exact if its genus
$\mathfrak g(A_i)=\log|\Delta_{A_i}|^{1/2}$ tends to infinity and the limits
$\phi_q=\lim N_q(A_i)/\mathfrak g(A_i)$,
$\phi_\R=\lim r_1(A_i)/\mathfrak g(A_i)$ and
$\phi_\C=\lim r_2(A_i)/\mathfrak g(A_i)$ exist. Their unconditional Basic
Inequality$'$ \cite[Section~3.2, Proposition~3.1]{TsfasmanVladut2002} states
that every such sequence satisfies
\[
 2\sum_q\phi_q\log q\sum_{m\ge1}\frac1{q^m+1}
 +\phi_\R\Bigl(\frac\gamma2+\frac12+\log2\sqrt\pi\Bigr)
 +\phi_\C(\gamma+\log4\pi)\le1.
\]
For our sequence $\mathfrak g(K^{(j)})=d_j\ell$ and $K^{(j)}$ is totally
imaginary, so $\phi_q=2\beta_q/\ell$, $\phi_\R=0$ and $\phi_\C=1/\ell$, and
the inequality becomes
\begin{equation}\label{an:tv-budget}
 \sum_q\beta_qw(q)\le\frac{\ell-\gamma-\log(4\pi)}2=2\kappa_\infty.
\end{equation}

(c) Since $g_1(q)\le1/(q-1)$ and $w(q)\ge2\log q/(q+1)$, we have
$g_1(q)\le3w(q)/(2\log2)$, so $Z(1)<\infty$ by (b). By
Lemma~\ref{an:shift} with $s-1$ in place of $\varepsilon$, and by (a), for
$1<s\le2$,
\[
 C'\le\limsup_j\frac{\log L_{F^{(j)}}(1)}{d_j}
 \le Z(s)+\sup_{0\le\theta\le1/2}\Upsilon(s-1,\theta).
\]
As $s\downarrow1$, $Z(s)$ increases to $Z(1)$ by monotone convergence, and
the supremum tends to zero. Hence $C'\le Z(1)$.
\end{proof}

The next proposition is the form of this argument with one abscissa. Here
it serves only for comparison (Remark~\ref{an:base-ceiling} and
Section~\ref{lim:section}) and for the Lean formalization
(Remark~\ref{an:lean-hypothesis}).

\begin{proposition}\label{an:prime-budget}
Let $\sigma=1+\varepsilon$ with $0<\varepsilon\le1$, and let $Y_*$ be a real
number with $[K:\Q]^{-1}\log\zeta_K(\sigma)\le Y_*$ for every
$K\in\mathcal K$. For each prime $\mathfrak p$ of $B$, let
$e^0_{\mathfrak p},f^0_{\mathfrak p}\ge1$ be integers with
$e_K(\mathfrak p)\ge e^0_{\mathfrak p}$ and $f_K(\mathfrak p)\ge
f^0_{\mathfrak p}$ for every $K\in\mathcal K$, and let $Y'_*$ be a real
number with
\[
 Y'_*\ge\sum_{\mathfrak p}\frac{\log\mathrm N\mathfrak p}
 {2e^0_{\mathfrak p}(\mathrm N\mathfrak p^{2f^0_{\mathfrak p}}-1)}.
\]
Then every $C'>Y_*+\varepsilon(\kappa_\infty+Y'_*)$ has the following
property: there is an $m_0$ such that $d^{-1}\log L_F(1)<C'$ for every
$K\in\mathcal K$ with $[K:B]\ge m_0$.
\end{proposition}

\begin{proof}
Suppose not, and let $K^{(j)}$, $\beta_q$ and $Z$ be as in
Lemma~\ref{an:limit}. For every prime power $q$,
\begin{align*}
 g_1(q)-g_{1+\varepsilon}(q)&=\int_1^{1+\varepsilon}\frac{\log q}{q^u-1}\,du
 \le\frac{\varepsilon\log q}{q-1},\\
 \frac{\log q}{q-1}-\frac{w(q)}2&=\log q\sum_{m\ge1}\frac1{q^m(q^m+1)}
 \le\frac{\log q}{q^2-1}.
\end{align*}
With \eqref{an:tv-budget} these give
$Z(1)\le Z(1+\varepsilon)+\varepsilon(\kappa_\infty+\sum_q\beta_q\log q/(q^2-1))$.
By Fatou's lemma, \eqref{an:local-debit} for $M=K^{(j)}$, and the monotonicity
in Lemma~\ref{an:local-contribution},
\[
 \sum_q\beta_q\frac{\log q}{q^2-1}
 \le\liminf_j\sum_{\mathfrak p}\frac{\log\mathrm N\mathfrak p}
 {2e_{K^{(j)}}(\mathfrak p)(\mathrm N\mathfrak p^{2f_{K^{(j)}}(\mathfrak p)}-1)}
 \le Y'_*.
\]
Finally $Z(1+\varepsilon)\le Y_*$ by Lemma~\ref{an:limit}(a). Thus
$C'\le Z(1)\le Y_*+\varepsilon(\kappa_\infty+Y'_*)$, a contradiction.
\end{proof}

By \eqref{an:local-debit} and the monotonicity in
Lemma~\ref{an:local-contribution}, $Y'_*$ bounds
$-[K:\Q]^{-1}(\zeta_K'/\zeta_K)(2)$ for every $K\in\mathcal K$, as $Y_*$
bounds $[K:\Q]^{-1}\log\zeta_K(\sigma)$.


\subsection{Selected Primes and Floors}\label{an:selected}
This subsection separates the primes of $B$ whose type is the same in every
$K\in\mathcal K$, the selected primes, whose contribution $Z_{\rm sel}(1)$
to $Z(1)$ is computed exactly (Lemma~\ref{an:selected-lemma}), from the free
primes, which enter Proposition~\ref{an:signed} only through lower bounds
for their residue degrees, the floors of Definition~\ref{an:floor-def}.

\begin{definition}\label{an:selected-def}
The \emph{selected primes} of $B$ are the six primes of $S$, which lie above
$2$, $3$ and $5$ (Lemma~\ref{tw:base}(c),(d)), and the capped primes
$\mathfrak t_1,\mathfrak t_2,\mathfrak t_3$
(Definition~\ref{tw:local-generators}). The other primes of $B$ are
\emph{free}. For a prime power $q$ put
\[
 \beta^{\rm sel}_q=\sum_{\mathfrak p}\frac1{2e_K(\mathfrak p)f_K(\mathfrak p)},
\]
the sum over the selected primes $\mathfrak p$ with
$\mathrm N\mathfrak p^{f_K(\mathfrak p)}=q$, where
$(e_K(\mathfrak p),f_K(\mathfrak p))$ is the type of $\mathfrak p$ in $K/B$
given by Theorem~\ref{tw:field-family}; it is the same for every
$K\in\mathcal K$. Put
\[
 Z_{\rm sel}(s)=\sum_q\beta^{\rm sel}_qg_s(q)\quad(s\ge1),\qquad
 T_{\rm sel}=\sum_q\beta^{\rm sel}_qw(q).
\]
The places of $F$ above the selected primes of $B$ form the set of finite
places used in the geometric construction (Section~\ref{sec:certificate});
in Subsection~\ref{intro:outline} they are also called selected.
\end{definition}

\begin{lemma}\label{an:selected-lemma}
The numbers $\beta^{\rm sel}_q$ are $1/32$ for $q=16$, $1/4$ for
$q=9,25,29^4$, $1/8$ for $q=41^4$, and $0$ for every other prime power $q$.
For every $K\in\mathcal K$ and every prime power $q$, the number
$N^{\rm sel}_q(K)$ of primes of $K$ of norm $q$ above selected primes of
$B$ satisfies $N^{\rm sel}_q(K)/[K:\Q]=\beta^{\rm sel}_q$. Moreover
\begin{gather*}
 0.0416684840321404<Z_{\rm sel}(1)<0.0416684840321405,\\
 0.2004555342790642<T_{\rm sel}<0.2004555342790643.
\end{gather*}
\end{lemma}

\begin{proof}
By Theorem~\ref{tw:field-family} the types are $(8,4)$ at
$\mathfrak p_1,\mathfrak p_2$, which have norm $2$; $(2,2)$ at the two primes
above $3$ and the two above $5$, which have norms $3$ and $5$; and $(1,4)$ at
$\mathfrak t_1,\mathfrak t_2$, of norm $29$, and at $\mathfrak t_3$, of norm
$41$. The definition then gives the values of $\beta^{\rm sel}_q$. A prime
$\mathfrak p$ of $B$ of type $(e,f)$ in the Galois extension $K/B$ has
$[K:B]/(ef)=d/(ef)$ primes of $K$ above it, each of norm
$\mathrm N\mathfrak p^f$, and $[K:\Q]=2d$; this gives
$N^{\rm sel}_q(K)/[K:\Q]=\beta^{\rm sel}_q$. The two numbers are evaluated
in Arb interval arithmetic by the supplementary program \texttt{signed.py}
(its functions \texttt{zsel} and \texttt{B\_r}).
\end{proof}

Since $K/B$ is unramified outside $S$, every free prime has $e_K(\mathfrak
p)=1$.

\begin{definition}\label{an:floor-def}
Put $q_{\min}=41^4=2825761$. A \emph{system of floors} is a choice of an
integer $f^0_{\mathfrak p}\ge1$, the \emph{floor} of $\mathfrak p$, for every
free prime $\mathfrak p$ of $B$ such that
\begin{enumerate}
\item[(Fl1)] $f^0_{\mathfrak p}$ divides $f_K(\mathfrak p)$ for every
 $K\in\mathcal K$;
\item[(Fl2)] $\mathrm N\mathfrak p^{f^0_{\mathfrak p}}\ge q_{\min}$.
\end{enumerate}
Its \emph{floor sum} is
\[
 Y_{\rm fl}(\sigma)=\sum_{\mathfrak p\ \rm free}
 \frac{g_\sigma(\mathrm N\mathfrak p^{f^0_{\mathfrak p}})}{2f^0_{\mathfrak p}}
 \qquad(\sigma>1).
\]
The series converges, since $g_\sigma(x^f)/f\le g_\sigma(x)$ for $x>1$ and
$f\ge1$, and $\sum_{\mathfrak p}g_\sigma(\mathrm N\mathfrak p)=\log\zeta_B(\sigma)$.
By (Fl1), the floors are an admissible choice, at the free primes, of the lower
bounds $f^0_{\mathfrak p}$ of Proposition~\fref{an:prime-budget}; this is why
the same symbol is used.
\end{definition}

By \eqref{an:local-sum}, the floor sum is the part of the termwise bound
for $[K:\Q]^{-1}\log\zeta_K(\sigma)$ that comes from the free primes, for a
hypothetical field in which every free prime has residue degree equal to
its floor. The floors used in this paper are those of
Definition~\ref{dh:floors}.

\subsection{A Signed Kernel}\label{an:signed-section}
Proposition~\ref{an:prime-budget} passes from $1+\varepsilon$ to $1$ through
the uniform bound $g_1(q)-g_{1+\varepsilon}(q)\le\varepsilon\log q/(q-1)$.
The next proposition, the criterion that proves Theorem~\ref{an:ceiling} in
Section~\ref{ce:section}, replaces it by a kernel inequality between $g_1$,
a combination $\mathsf K$ of several functions $g_{\sigma_i}$ with
coefficients of either sign, and the weight $w$ of the
Tsfasman--Vl\u{a}du\c{t} inequality. Such a combination can follow the
function $(g_1-b_{\rm TV}w)_+$, where $b_{\rm TV}\ge0$ is the multiplier of
$w$ in the proposition, far more closely than a single multiple of
$g_\sigma$.

\begin{proposition}\label{an:signed}
Let $(f^0_{\mathfrak p})$ be a system of floors (Definition~\ref{an:floor-def}),
with floor sum $Y_{\rm fl}$, and let $Z_{\rm sel}$ and $T_{\rm sel}$ be as in
Definition~\ref{an:selected-def} and $\kappa_\infty$ as in
Subsection~\ref{an:budget-section}.
Let $1<\sigma_1<\dots<\sigma_n$, let $c_1,\dots,c_n$ be real numbers of
either sign, let $b_{\rm TV}\ge0$, and define the \emph{kernel}
$\mathsf K(q)=\sum_ic_ig_{\sigma_i}(q)$ for real $q>1$. Suppose that, for
every real $q\ge q_{\min}$,
\[
 \text{(i) }\mathsf K(q)\ge g_1(q)-b_{\rm TV}\,w(q),\qquad
 \text{(ii) }\mathsf K(q)\ge0,\qquad
 \text{(iii) }\mathsf K\text{ is non-increasing on }[q_{\min},\infty).
\]
Then every
\[
 C'>Z_{\rm sel}(1)+\sum_{i=1}^nc_iY_{\rm fl}(\sigma_i)
 +b_{\rm TV}\,(2\kappa_\infty-T_{\rm sel})
\]
has the following property: there is an $m_0$ such that
$d^{-1}\log L_F(1)<C'$ for every $K\in\mathcal K$ with $[K:B]\ge m_0$.
\end{proposition}

\begin{proof}
Suppose not, and let $K^{(j)}$, $\beta_q$ and $Z$ be as in
Lemma~\ref{an:limit}, applied with this $C'$. Write
$N_q(K)=N^{\rm sel}_q(K)+N^{\rm free}_q(K)$, where $N^{\rm free}_q(K)$
counts the primes of $K$ of norm $q$ above free primes of $B$, and put
$\mu_j(q)=N^{\rm free}_q(K^{(j)})/(2d_j)$. By Lemma~\ref{an:selected-lemma},
$\beta^{\rm free}_q=\beta_q-\beta^{\rm sel}_q=\lim_j\mu_j(q)$ for every $q$,
and $\beta^{\rm free}_q\ge0$.

\emph{One field.} Fix $j$ and write $K=K^{(j)}$ and $d=d_j$. A free prime
$\mathfrak p$ of $B$ with $f=f_K(\mathfrak p)$ has $d/f$ primes of $K$ above
it, each of norm $\mathrm N\mathfrak p^f$. By (Fl1), $f$ is a multiple of
$f^0_{\mathfrak p}$, so $1/(2f)\le1/(2f^0_{\mathfrak p})$, and
$\mathrm N\mathfrak p^f\ge\mathrm N\mathfrak p^{f^0_{\mathfrak p}}\ge q_{\min}$
by (Fl2). By (ii) and (iii), the primes of $K$ above $\mathfrak p$ contribute
\[
 \frac{d/f}{2d}\,\mathsf K(\mathrm N\mathfrak p^f)
 =\frac{\mathsf K(\mathrm N\mathfrak p^f)}{2f}
 \le\frac{\mathsf K(\mathrm N\mathfrak p^{f^0_{\mathfrak p}})}{2f^0_{\mathfrak p}}
\]
to $\sum_q\mu_j(q)\mathsf K(q)$. Summing over the free primes,
\[
 \sum_q\mu_j(q)\mathsf K(q)\le\Sigma_{\mathsf K}:=\sum_{\mathfrak p\ \rm free}
 \frac{\mathsf K(\mathrm N\mathfrak p^{f^0_{\mathfrak p}})}{2f^0_{\mathfrak p}}
 =\sum_{i=1}^nc_iY_{\rm fl}(\sigma_i),
\]
where the series defining $\Sigma_{\mathsf K}$ converges absolutely, since
$|\mathsf K|\le\sum_i|c_i|g_{\sigma_i}$, and may be summed term by term. Note
also that $\mu_j$ is supported on $[q_{\min},\infty)$, hence so is
$\beta^{\rm free}$.

\emph{The limit.} Since $\mathsf K\ge0$ on $[q_{\min},\infty)$, Fatou's
lemma gives $\sum_q\beta^{\rm free}_q\mathsf K(q)\le\liminf_j\sum_q\mu_j(q)
\mathsf K(q)\le\Sigma_{\mathsf K}$. By (i) at every $q$ in the support of
$\beta^{\rm free}$, by $b_{\rm TV}\ge0$, and by Lemma~\ref{an:limit}(b),
\[
 \sum_q\beta^{\rm free}_qg_1(q)
 \le\sum_q\beta^{\rm free}_q\mathsf K(q)
 +b_{\rm TV}\sum_q\beta^{\rm free}_qw(q)
 \le\Sigma_{\mathsf K}+b_{\rm TV}\,(2\kappa_\infty-T_{\rm sel}),
\]
where we used $\sum_q\beta^{\rm free}_qw(q)=\sum_q\beta_qw(q)-T_{\rm sel}$.
With Lemma~\ref{an:limit}(c),
\[
 C'\le Z(1)=Z_{\rm sel}(1)+\sum_q\beta^{\rm free}_qg_1(q)
 \le Z_{\rm sel}(1)+\sum_ic_iY_{\rm fl}(\sigma_i)
 +b_{\rm TV}\,(2\kappa_\infty-T_{\rm sel}),
\]
which contradicts the choice of $C'$.
\end{proof}

\begin{remark}\label{an:signed-remark}
Condition (ii) follows from (iii), since $\mathsf K(q)\to0$ as
$q\to\infty$; it is checked separately. The proof uses the floors only
through (Fl1) and (Fl2), and the selected primes only through their exact
types. The right-hand side of the proposition is at least
$Z_{\rm sel}(1)$, since $\Sigma_{\mathsf K}\ge0$ by (ii) and
$2\kappa_\infty>T_{\rm sel}$ (Remark~\ref{an:limit-remark}). With $n=1$ and
$c_1\ge1$ it is the refined Tsfasman--Vl\u{a}du\c{t} step of
Proposition~\ref{lim:D}, without the adaptive credits
$R_{\mathfrak p}(f_K(\mathfrak p))+\rho_K(\mathfrak p)I_{\mathfrak p}(f_K(\mathfrak p))$
of Definition~\ref{lim:credits} and Proposition~\ref{lim:A}.
\end{remark}

\subsection{An Approximate Functional Equation for Two Gamma Factors}
\label{an:afe-section}
This subsection and the next two enclose $\zeta_E(\sigma)$ at the abscissae
of Section~\ref{ce:section} (Proposition~\ref{an:genus-value}), through the
factorization of Proposition~\ref{gf:factorization}. The factors
$L(\sigma,\chi_e)$ are computed from an approximate functional equation
(Lemma~\ref{an:balanced-afe}), whose kernels are bounded by envelopes
(Lemmas~\ref{an:kernels} and~\ref{an:tails}) and enclosed on a grid
(Subsection~\ref{an:kernel-enclosures}). For real $\mu$ and $x,y>0$ define
the incomplete gamma and Bessel integrals
\[
 H_\mu(x)=\int_1^\infty u^{\mu-1}e^{-xu}\,du,\qquad
 I_\mu(y)=\int_1^\infty u^\mu K_0(yu)\,du,
\]
where $K_0$ is the modified Bessel function of the second kind of order
zero, and their envelopes (upper bounds, by Lemma~\ref{an:kernels}(2))
\[
 \mathcal V_1(x)=\frac{e^{-x}}x\Bigl(1+\frac1x\Bigr),\qquad
 \mathcal V_2(y)=\sqrt{\frac\pi2}\,\frac{e^{-y}}{y^{3/2}}\Bigl(1+\frac1y\Bigr).
\]
Numerical evaluation of $L$-functions from their functional equations is
developed in \cite{Dokchitser2004}. Part (1) of the next lemma treats mixed
gamma factors, and part (2), for pure gamma factors, uses the Bessel
function $K_0$. We prove both.

\begin{lemma}\label{an:balanced-afe}
Let $a_n$ be real numbers with $|a_n|\le\tau(n)$, where $\tau(n)$ is the
number of divisors of $n$, let $Q>0$ and
$\nu_1,\nu_2\in\{0,1\}$ (integers, not places), and let
$L(s)=\sum_{n\ge1}a_nn^{-s}$. Suppose that
$\Lambda(s)=Q^{s/2}\Gamma_\R(s+\nu_1)\Gamma_\R(s+\nu_2)L(s)$ extends to an
entire function of order at most one with $\Lambda(1-s)=\Lambda(s)$. Put
$t=2\pi/\sqrt Q$ and let $\sigma>1$ be real.
\begin{enumerate}
\item If $\nu_1\ne\nu_2$, then
\[
 L(\sigma)=\frac{t^\sigma}{\Gamma(\sigma)}\sum_{n\ge1}a_n
 \bigl(H_\sigma(tn)+H_{1-\sigma}(tn)\bigr).
\]
\item If $\nu_1=\nu_2=\nu$, then
\[
 L(\sigma)=\frac{4(t/2)^{\sigma+\nu}}{\Gamma((\sigma+\nu)/2)^2}
 \sum_{n\ge1}a_nn^\nu
 \bigl(I_{\sigma+\nu-1}(tn)+I_{\nu-\sigma}(tn)\bigr).
\]
\end{enumerate}
The series converge absolutely.
\end{lemma}

\begin{proof}
In case (1) put $k(x)=e^{-x}$ and $\nu=0$, and in case (2) put
$k(x)=K_0(x)$. Their Mellin transforms are $\Gamma(w)$ and
$2^{w-2}\Gamma(w/2)^2$ for $\operatorname{Re}w>0$
\cite[10.43.19]{DLMF}. Let $\Theta(u)=\sum_na_nn^\nu k(tnu)$ for $u>0$. For
$\operatorname{Re}w>1+\nu$, termwise integration gives
\[
 \widetilde\Theta(w)=\int_0^\infty\Theta(u)u^{w-1}\,du=
 \begin{cases}
 t^{-w}\Gamma(w)L(w)=\tfrac12\Lambda(w)&\text{in case (1)},\\
 2^{w-2}t^{-w}\Gamma(w/2)^2L(w-\nu)=\tfrac14Q^{\nu/2}\Lambda(w-\nu)
 &\text{in case (2)}.
 \end{cases}
\]
In case (1) we used $\Gamma_\R(w)\Gamma_\R(w+1)=\Gamma_\C(w)$ and
$Q^{w/2}(2\pi)^{-w}=t^{-w}$. In both cases $\widetilde\Theta$ is entire and
$\widetilde\Theta(w)=\widetilde\Theta(1+2\nu-w)$.

The function $L=\Lambda/(Q^{s/2}\Gamma_\R(s+\nu_1)\Gamma_\R(s+\nu_2))$ of the
complex variable $s$ is entire of finite order. Fix $\eta>0$. The function
$L$ is bounded on $\operatorname{Re}s=1+\eta$, and, by the functional
equation and Stirling's formula, of polynomial growth on
$\operatorname{Re}s=-\eta$. By the Phragm\'en--Lindel\"of principle it has
polynomial growth in $|\operatorname{Im}s|$ on every vertical strip. With
Stirling's formula this shows that $\widetilde\Theta(w)$ decays exponentially in
$|\operatorname{Im}w|$, uniformly on vertical strips. Mellin inversion on a
line $\operatorname{Re}w=c>1+\nu$ and a shift to
$\operatorname{Re}w=1+2\nu-c$ are therefore justified, and no residues
occur. With the functional equation of $\widetilde\Theta$ they give
\[
 \Theta(u)=u^{-1-2\nu}\,\Theta(1/u)\qquad(u>0).
\]
In particular $\Theta(u)$ decays rapidly as $u\downarrow0$. Splitting the
Mellin integral at $u=1$ and substituting $u\mapsto1/u$ below $1$ gives
\[
 \widetilde\Theta(\sigma+\nu)=\int_1^\infty\Theta(u)
 \bigl(u^{\sigma+\nu-1}+u^{\nu-\sigma}\bigr)\,du.
\]
Termwise integration, justified by the exponential decay of $k$, turns the
right side into $\sum_na_nn^\nu$ times $H_\sigma(tn)+H_{1-\sigma}(tn)$ in case
(1), and times $I_{\sigma+\nu-1}(tn)+I_{\nu-\sigma}(tn)$ in case (2). On the
left, $\widetilde\Theta(\sigma)=t^{-\sigma}\Gamma(\sigma)L(\sigma)$ in case (1), and
$\widetilde\Theta(\sigma+\nu)=2^{\sigma+\nu-2}t^{-\sigma-\nu}
\Gamma((\sigma+\nu)/2)^2L(\sigma)$ in case (2).
\end{proof}

\begin{lemma}\label{an:kernels}
Let $\mu\in\R$.
\begin{enumerate}
\item The functions $H_\mu$ and $I_\mu$ are Laplace transforms of positive
 measures on $[1,\infty)$. They are completely monotone on $(0,\infty)$,
 extend holomorphically to $\operatorname{Re}z>0$, and satisfy
 $|H(z)|\le H(\operatorname{Re}z)$ there, for $H=H_\mu$ or $H=I_\mu$.
\item If $\mu\le2$, then $H_\mu\le\mathcal V_1$. If $\mu\le3/2$, then
 $I_\mu\le\mathcal V_2$.
\item They satisfy
 \[
 xH_\mu'(x)+\mu H_\mu(x)=-e^{-x},\qquad
 yI_\mu'(y)+(\mu+1)I_\mu(y)=-K_0(y),
 \]
 and $K_0$ satisfies $y^2K_0''+yK_0'-y^2K_0=0$ and $K_0'=-K_1$, where
 $K_1$ is the modified Bessel function of the second kind of order one.
\end{enumerate}
\end{lemma}

\begin{proof}
The definition exhibits $H_\mu$ as the Laplace transform of
$u^{\mu-1}\,du$ on $[1,\infty)$. Taking the order $0$ in
\cite[10.32.9]{DLMF} and substituting $v=\cosh r$ in its integral gives
$K_0(x)=\int_1^\infty e^{-xv}(v^2-1)^{-1/2}\,dv$. Substituting this into
$I_\mu$, and then $v=r/u$ in the inner integral, shows
\[
 I_\mu(y)=\int_1^\infty e^{-yr}\Bigl(\int_1^ru^\mu(r^2-u^2)^{-1/2}\,du\Bigr)
 \,dr,
\]
with a nonnegative inner integral. Differentiation under the integral
gives complete monotonicity. Since the measures are positive and
$|e^{-zr}|=e^{-r\operatorname{Re}z}$, we get $|H(z)|\le H(\operatorname{Re}z)$
for $\operatorname{Re}z>0$, which proves (1).

For (2), $u^{\mu-1}\le u$ for $u\ge1$ and $\mu\le2$, and
$\int_1^\infty ue^{-xu}\,du=\mathcal V_1(x)$. For $I_\mu$ we use
$K_0(x)\le K_{1/2}(x)=\sqrt{\pi/(2x)}\,e^{-x}$ \cite[10.39.2]{DLMF}. It
holds because the modified Bessel function of order $\eta\ge0$ is
$\int_0^\infty e^{-x\cosh r}\cosh(\eta r)\,dr$ \cite[10.32.9]{DLMF}, which
increases with $\eta$. Since $u^{\mu-1/2}\le u$ for $\mu\le3/2$,
\[
 I_\mu(y)\le\sqrt{\frac\pi{2y}}\int_1^\infty u^{\mu-1/2}e^{-yu}\,du
 \le\mathcal V_2(y).
\]

For (3), differentiate under the integral and integrate by parts:
$xH_\mu'(x)=-\int_1^\infty u^\mu xe^{-xu}\,du=-e^{-x}-\mu H_\mu(x)$, and
$yI_\mu'(y)=\int_1^\infty u^{\mu+1}\frac{d}{du}K_0(yu)\,du
=-K_0(y)-(\mu+1)I_\mu(y)$. The last two identities are Bessel's equation
and a standard derivative \cite[10.25.1, 10.29.3]{DLMF}.
\end{proof}

The next lemma bounds the tails of the two series of
Lemma~\ref{an:balanced-afe}. In the application $1<\sigma\le101/100$, so its
hypothesis $1<\sigma\le3/2$ holds.


\begin{lemma}\label{an:tails}
Let $1<\sigma\le3/2$, let $|a_n|\le\tau(n)$, let $t>0$, $N\ge1$ and
$\omega>1$.
\begin{enumerate}
\item If $t(N+1)\ge\omega-1$, then for $\mu\in\{\sigma,1-\sigma\}$
\[
 \sum_{n>N}|a_n|H_\mu(tn)\le\zeta(\omega)^2(N+1)^\omega
 \mathcal V_1\bigl(t(N+1)\bigr).
\]
\item If $\nu\in\{0,1\}$ and $t(N+1)\ge\omega+\nu-\frac32$, then for
 $\mu\in\{\sigma+\nu-1,\nu-\sigma\}$
\[
 \sum_{n>N}|a_n|n^\nu I_\mu(tn)\le\zeta(\omega)^2(N+1)^{\omega+\nu}
 \mathcal V_2\bigl(t(N+1)\bigr).
\]
\end{enumerate}
\end{lemma}

\begin{proof}
Since $1<\sigma\le3/2$ and $\nu\in\{0,1\}$, every $\mu\in\{\sigma,1-\sigma\}$
satisfies $\mu\le2$ and every $\mu\in\{\sigma+\nu-1,\nu-\sigma\}$ satisfies
$\mu\le3/2$; so Lemma~\ref{an:kernels}(2) gives $H_\mu\le\mathcal V_1$ in (1)
and $I_\mu\le\mathcal V_2$ in (2). The logarithmic derivative of
$x^\omega\mathcal V_1(x)=x^{\omega-1}e^{-x}(1+1/x)$ is
$(\omega-1)/x-1-1/(x(x+1))<0$ for $x\ge\omega-1$. So
$n^\omega\mathcal V_1(tn)\le(N+1)^\omega\mathcal V_1(t(N+1))$ for $n>N$.
Hence
\[
 \sum_{n>N}|a_n|H_\mu(tn)\le\sum_{n>N}\frac{\tau(n)}{n^\omega}\,
 n^\omega\mathcal V_1(tn)\le(N+1)^\omega\mathcal V_1\bigl(t(N+1)\bigr)
 \sum_{n\ge1}\frac{\tau(n)}{n^\omega},
\]
and the last sum is $\zeta(\omega)^2$. Part (2) is the same, since
$y^{\omega+\nu}\mathcal V_2(y)$ is a multiple of
$y^{\omega+\nu-3/2}e^{-y}(1+1/y)$, which decreases for
$y\ge\omega+\nu-\frac32$.
\end{proof}

\subsection{\texorpdfstring{Enclosures of $H_\mu$ and $I_\mu$, and Finite Sums}{Enclosures of H and I, and Finite Sums}}
\label{an:kernel-enclosures}
This subsection encloses the kernels $H_\mu$ and $I_\mu$ on a geometric grid
(Lemma~\ref{an:grid}) and bounds the error of the resulting approximation of
the finite sums of Lemma~\ref{an:balanced-afe} (Lemma~\ref{an:binning});
Lemma~\ref{an:row} uses the second, which rests on the first. An \emph{enclosure} of a real number is an
interval that contains it. Let $H$ be $H_\mu$ or $I_\mu$
with $\mu$ as in Lemma~\ref{an:tails}, and let $\mathcal V$ be its envelope,
$\mathcal V_1$ or $\mathcal V_2$. For $J\ge0$ and $\xi>0$ put
\[
 \rho_J(\xi)=\mathcal V(\xi/2)\,\frac{16^{-J-1}}{1-1/16}.
\]

\begin{lemma}\label{an:taylor}
Let $\xi>0$ and let $h_j=H^{(j)}(\xi)/j!$ be the Taylor coefficients of $H$
at $\xi$.
\begin{enumerate}
\item $(-1)^jh_j\ge0$ and $|h_j|\le\mathcal V(\xi/2)(2/\xi)^j$ for all
 $j\ge0$.
\item For $|x-\xi|\le\xi/32$,
 $\bigl|H(x)-\sum_{j=0}^Jh_j(x-\xi)^j\bigr|\le\rho_J(\xi)$. For
 $x=31\xi/32$ the difference lies in $[0,\rho_J(\xi)]$.
\item Let $k_j$ be the Taylor coefficients at $\xi$ of $e^{-x}$ if
 $H=H_\mu$, and of $K_0$ if $H=I_\mu$. If $H=H_\mu$, then
 $(j+1)\xi h_{j+1}=-k_j-(j+\mu)h_j$, and $k_j=(-1)^je^{-\xi}/j!$. If
 $H=I_\mu$, then $(j+1)\xi h_{j+1}=-k_j-(j+\mu+1)h_j$, where
 $k_0=K_0(\xi)$, $k_1=-K_1(\xi)$, and
 \[
 (j+1)(j+2)\xi^2k_{j+2}=-(j+1)(2j+1)\xi k_{j+1}-(j^2-\xi^2)k_j
 +2\xi k_{j-1}+k_{j-2},
 \]
 with $k_{-1}=k_{-2}=0$.
\end{enumerate}
\end{lemma}

\begin{proof}
Complete monotonicity gives the signs in (1). The disc $|z-\xi|\le\xi/2$ lies
in $\operatorname{Re}z\ge\xi/2$, where $|H(z)|\le H(\xi/2)\le\mathcal
V(\xi/2)$ by Lemma~\ref{an:kernels}. Cauchy's estimate gives the bound on
$h_j$. For $|x-\xi|\le\xi/32$, the omitted terms have modulus at most
$\mathcal V(\xi/2)16^{-j}$, and their sum over $j>J$ is at most
$\rho_J(\xi)$. For $x<\xi$ every term $h_j(x-\xi)^j$ is nonnegative by (1),
so the remainder is nonnegative. Part (3) follows by comparing coefficients
of $(x-\xi)^j$ in the differential equations of Lemma~\ref{an:kernels}(3),
written at $x=\xi+(x-\xi)$.
\end{proof}

The grid points are $\xi_i=64\,(31/32)^i$ for $0\le i\le i_1$, where
$i_1=1878$ is the least index with $\xi_{i_1}\le2^{-80}$. Taylor polynomials
are formed at $\xi_0,\ldots,\xi_{i_1-1}$. The initial values
$e^{-\xi_i}$, $K_0(\xi_i)$ and $K_1(\xi_i)$ of part~(3) of
Lemma~\ref{an:taylor} are enclosed with the rigorous exponential and Bessel
functions of Arb \cite{Johansson2017}, which is now part of FLINT
\cite{FLINT} and is called through its Python interface.

\begin{lemma}\label{an:grid}
Let $J\ge0$, and suppose that enclosures of $e^{-\xi_i}$, $K_0(\xi_i)$ and
$K_1(\xi_i)$ are given for $0\le i<i_1$. Start from the enclosure
$[0,\mathcal V(64)]$ of $H(\xi_0)$. For $i=0,1,\ldots,i_1-1$, compute from
the enclosure of $H(\xi_i)$ enclosures of $h_0,\ldots,h_J$ at $\xi_i$ by the
recurrences of Lemma~\ref{an:taylor}(3), and take
\[
 \sum_{j=0}^Jh_j(-\xi_i/32)^j+[0,\rho_J(\xi_i)],
\]
evaluated in interval arithmetic, as the enclosure of $H(\xi_{i+1})$. Then
every interval so obtained contains the corresponding true value.
\end{lemma}

\begin{proof}
By Lemma~\ref{an:kernels}, $H(\xi_0)\in[0,\mathcal V(64)]$. Suppose that
the enclosure of $H(\xi_i)$ contains $H(\xi_i)$. The true Taylor
coefficients at $\xi_i$ are obtained from the true value $H(\xi_i)$ by the
same recurrence, so their enclosures contain them. Since
$\xi_{i+1}=31\xi_i/32$, Lemma~\ref{an:taylor}(2) shows that the enclosure
of $H(\xi_{i+1})$ contains $H(\xi_{i+1})$. Induction on $i$ proves the
lemma.
\end{proof}

Now let $t>0$, $\nu\in\{0,1\}$, integers $a_1,\ldots,a_N$ and $J\ge0$ be
given, with $tN<\xi_0$ and $t>\xi_{i_1}$. For $0\le i<i_1$ let
$\mathrm{Bin}_i=\{n:1\le n\le N,\ \xi_{i+1}<tn\le\xi_i\}$. These sets
partition $\{1,\ldots,N\}$. For nonempty $\mathrm{Bin}_i$, let $\bar n_i$ be
the integer part of the average of its least and largest elements, and put
\[
 M_{i,j}=\sum_{n\in\mathrm{Bin}_i}a_nn^\nu(n-\bar n_i)^j,\qquad
 A_i=\sum_{n\in\mathrm{Bin}_i}|a_n|n^\nu,\qquad
 P_{i,j}=t^j\sum_{m=j}^J\binom mjh_{i,m}(t\bar n_i-\xi_i)^{m-j},
\]
where $h_{i,m}$ are the Taylor coefficients of $H$ at $\xi_i$. The moments
$M_{i,j}$ and the absolute sums $A_i$ are exact integers.

\begin{lemma}\label{an:binning}
With this notation,
\begin{equation}\label{an:bin-sum}
 \Bigl|\sum_{n=1}^Na_nn^\nu H(tn)-\sum_i\sum_{j=0}^JP_{i,j}M_{i,j}\Bigr|
 \le\sum_i\rho_J(\xi_i)A_i,
\end{equation}
where $i$ runs over the indices with $\mathrm{Bin}_i$ nonempty.
\end{lemma}

\begin{proof}
Let $n\in\mathrm{Bin}_i$. Then $|tn-\xi_i|<\xi_i/32$, so by
Lemma~\ref{an:taylor}(2) the value $H(tn)$ differs from
$\sum_{m=0}^Jh_{i,m}(tn-\xi_i)^m$ by at most $\rho_J(\xi_i)$. Since
$tn-\xi_i=t(n-\bar n_i)+(t\bar n_i-\xi_i)$, the binomial theorem gives
$\sum_{m=0}^Jh_{i,m}(tn-\xi_i)^m=\sum_{j=0}^JP_{i,j}(n-\bar n_i)^j$.
Multiply by $a_nn^\nu$ and sum over $n\in\mathrm{Bin}_i$ and over $i$.
\end{proof}

In the application below $t\ge2\pi/\sqrt{3470400}>0.0033$, because
$Q_e\le3470400$ for every $e$ (Corollary~\ref{gf:conductors}). So every argument
$tn$ exceeds $0.0033$, and the part of the grid below that point is not
used.

\subsection{The Values of $\zeta_E$}
\label{an:value-section}
This subsection encloses $Y_E(\sigma)=\log\zeta_E(\sigma)/512$ at the
abscissae of Table~\ref{ce:kernel-table} (Proposition~\ref{an:genus-value});
these values enter $Y_{W''}$ through \eqref{dh:YW}.

Let $\sigma\in(1,101/100]$ be rational, and let $e\in\Ftwo^8\setminus\{0\}$.
By Proposition~\ref{gf:degree-two}, $L(s,\chi_e)=\sum_na_n(e)n^{-s}$
satisfies the hypotheses of Lemma~\ref{an:balanced-afe} with $Q=Q_e$ and
$\nu_k=\nu_k(e)$ (Definition~\ref{gf:archimedean}). Put $t=2\pi/\sqrt{Q_e}$ and let the truncation point be
\[
 N^{\rm afe}_e=8\lfloor\sqrt{Q_e}\rfloor+1,
\]
twice the length $4\lfloor\sqrt{Q_e}\rfloor+1$ of Remark~\ref{gf:checks},
so that the enclosures are narrow enough for Proposition~\ref{an:signed}.

\begin{lemma}\label{an:afe-range}
For every $e\in\Ftwo^8\setminus\{0\}$, with $t=2\pi/\sqrt{Q_e}$ and
$N=N^{\rm afe}_e$,
\[
 tN<50.6<\xi_0,\qquad t>0.0033>\xi_{i_1},\qquad 49.0<t(N+1)<50.6.
\]
Hence Lemma~\ref{an:binning} applies with this $t$ and $N$, and for
$\nu\in\{0,1\}$ the conditions on $t(N+1)$ in Lemma~\ref{an:tails} hold for every
$\omega$ in the list
$\frac{101}{100},\frac{21}{20},\frac{11}{10},\frac98,\frac65,\frac54,
\frac43,\frac32,2$.
\end{lemma}

\begin{proof}
By Corollary~\ref{gf:conductors}, $723\le Q_e\le3470400$. Since
$\sqrt{Q_e}-1<\lfloor\sqrt{Q_e}\rfloor\le\sqrt{Q_e}$,
\[
 tN\le2\pi\Bigl(8+\frac1{\sqrt{Q_e}}\Bigr)\le2\pi\Bigl(8+\frac1{\sqrt{723}}\Bigr)
 <50.6,\qquad
 t(N+1)>2\pi\Bigl(8-\frac6{\sqrt{Q_e}}\Bigr)\ge2\pi\Bigl(8-\frac6{\sqrt{723}}\Bigr)
 >48.8,
\]
and $t\ge2\pi/\sqrt{3470400}>0.0033$. Moreover $\xi_0=64$ and
$\xi_{i_1}\le2^{-80}$. The conditions on $t(N+1)$ in Lemma~\ref{an:tails} are
$t(N+1)\ge\omega-1$ and $t(N+1)\ge\omega+\nu-\frac32$, whose right sides are
at most $\frac32$ for $\omega\le2$ and $\nu\le1$. Finally, $Q_e$ takes the
$53$ values $241\cdot2^i\cdot m$ with a nonzero entry in
\eqref{gf:conductor-table}, and evaluating $t(N+1)$ at these values gives
$49.0<t(N+1)<50.6$; the extremes, $49.07$ and $50.59$, occur at
$Q_e=723$ and $Q_e=964$.
\end{proof}

Let $\Pi_e$ be the prefactor of Lemma~\ref{an:balanced-afe}:
$\Pi_e=t^\sigma/\Gamma(\sigma)$ if $\nu_1(e)\ne\nu_2(e)$, and
$\Pi_e=4(t/2)^{\sigma+\nu}/\Gamma((\sigma+\nu)/2)^2$ if
$\nu_1(e)=\nu_2(e)=\nu$; in the first case put $\nu=0$. Let $S^{(1)}$ and
$S^{(2)}$ be the sums $\sum_i\sum_{j=0}^{20}P_{i,j}M_{i,j}$ of
Lemma~\ref{an:binning}, with $N=N^{\rm afe}_e$, $a_n=a_n(e)$ and $J=20$, for
the two functions of the series of Lemma~\ref{an:balanced-afe}: $H=H_\sigma$
and $H=H_{1-\sigma}$ in the first case, and $H=I_{\sigma+\nu-1}$ and
$H=I_{\nu-\sigma}$ in the second; their envelope is $\mathcal V=\mathcal V_1$
in the first case and $\mathcal V=\mathcal V_2$ in the second. By
Lemma~\ref{an:afe-range}, Lemma~\ref{an:binning} applies, and
Lemma~\ref{an:tails} applies with every $\omega$ in the list above; the value
of $\omega$ in this list is chosen to minimize the bound
\[
 \mathcal T_e=\sum_i\rho_{20}(\xi_i)A_i
 +\zeta(\omega)^2(N^{\rm afe}_e+1)^{\omega+\nu}
 \mathcal V\bigl(t(N^{\rm afe}_e+1)\bigr).
\]

\begin{lemma}\label{an:row}
For rational $\sigma\in(1,101/100]$ and every $e\in\Ftwo^8\setminus\{0\}$,
\begin{equation}\label{an:row-enclosure}
 \bigl|L(\sigma,\chi_e)-\Pi_e(S^{(1)}+S^{(2)})\bigr|
 \le2\Pi_e\mathcal T_e.
\end{equation}
\end{lemma}

\begin{proof}
Split each of the two series of Lemma~\ref{an:balanced-afe} at
$N^{\rm afe}_e$. By Lemma~\ref{an:binning}, each finite part differs from its
approximation $S^{(1)}$ or $S^{(2)}$ by at most $\sum_i\rho_{20}(\xi_i)A_i$; the
same bound serves both series, since their two functions have the same
envelope. By Lemma~\ref{an:tails}, each tail is at most
$\zeta(\omega)^2(N^{\rm afe}_e+1)^{\omega+\nu}\mathcal V(t(N^{\rm afe}_e+1))$. Multiplying by
$\Pi_e$ gives \eqref{an:row-enclosure}.
\end{proof}

\begin{proposition}\label{an:genus-value}
Put $Y_E(\sigma)=\log\zeta_E(\sigma)/512$. For each of the abscissae
$\sigma$ of Table~\ref{ce:kernel-table}, the supplementary program
\texttt{ye\_hp.py} computes an enclosure of $Y_E(\sigma)$ of width less than
$1.3\cdot10^{-19}$, by \eqref{an:row-enclosure} and
\[
 \log\zeta_E(\sigma)=\log\zeta(\sigma)+\log L(\sigma,\chi_B)
 +\sum_{e\ne0}\log L(\sigma,\chi_e).
\]
For example,
\[
 0.085358048288873992<Y_E(1001/1000)<0.085358048288873993.
\]
\end{proposition}

\begin{proof}[Proof (computer-assisted)]
The factorization is Proposition~\ref{gf:factorization}. The coefficients
$a_n(e)$ for $n\le N^{\rm afe}_e$ are computed exactly from the Euler factors of
Proposition~\ref{gf:local-data} by the functions of the supplementary program
\texttt{lfun241.py}, and the program \texttt{afe241.py}, both in the
directory \path{papers/0.04273/certificates} of \cite{NaslundCode4317} and
called unchanged by \texttt{ye\_hp.py}, evaluates
\eqref{an:row-enclosure} in Arb midpoint--radius interval arithmetic
\cite{Johansson2017,FLINT} with $256$ bits of precision, rounding every
quantity outward. Its routines for
Lemmas~\ref{an:balanced-afe}--\ref{an:binning} are taken unchanged from the
data archive in \path{papers/0.0418235/certificates} of
\cite{NaslundCode4317}, whose members are checked against their SHA-256
hashes when it is unpacked, and they implement these lemmas as stated here. Every enclosure of
$L(\sigma,\chi_e)$ is positive. The factor $\zeta_B(\sigma)=\zeta(\sigma)
L(\sigma,\chi_B)$ is evaluated with
\[
 L(\sigma,\chi_B)=241^{-\sigma}\sum_{a=1}^{240}\Bigl(\frac
 a{241}\Bigr)\zeta\Bigl(\sigma,\frac a{241}\Bigr),
\]
where $\zeta(s,x)=\sum_{n\ge0}(n+x)^{-s}$ is the Hurwitz zeta function,
using Arb's rigorous implementation of it \cite{Johansson2015Hurwitz}.
\end{proof}

\subsection{Frobenius Vectors and the Set $\mathcal P_4$}
\label{an:census}
This subsection defines a set $\mathcal P_4$ of free primes of norm at
most $10^6$ by a condition on their Frobenius vectors
(Definition~\ref{an:P4}), and shows that their residue degree in every
$K\in\mathcal K$ is divisible by $4$ (Lemma~\ref{an:census-degree}); these
primes receive the floor
$f^0_{\mathfrak p}=4$ in Definition~\ref{dh:floors}. We use the group $G_B$ of Definition~\ref{tw:GB} and its Zassenhaus
filtration $D_nG_B$, the relation space $R_2$ in degree two
(Definition~\ref{tw:relation-space}), which Lemma~\ref{tw:quadratic-layer}
computes, the restricted square $v^{[2]}$ of a vector $v\in\Ftwo^8$ as in
Section~\ref{tw:tower}, and the Frobenius vectors
$\overline{\Frob}_{\mathfrak p}$ of Definition~\ref{gf:frobenius-vector}.

\begin{definition}\label{an:P4}
Let $\Sigma_2=\{v\in\Ftwo^8:v^{[2]}\in R_2\}$, and let $\mathcal P_4$ be the
set of free primes $\mathfrak p$ of $B$ with $\mathrm N\mathfrak p\le10^6$
and $\overline{\Frob}_{\mathfrak p}\notin\Sigma_2$.
\end{definition}

\begin{lemma}\label{an:census-degree}
Let $K\in\mathcal K$, and let $\mathfrak p\notin S$ be a prime of $B$ with
$\overline{\Frob}_{\mathfrak p}\notin\Sigma_2$. Then $e_K(\mathfrak p)=1$, and
$f_K(\mathfrak p)$ is a power of two and at least $4$, so $4$ divides
$f_K(\mathfrak p)$.
\end{lemma}

\begin{proof}
The field $K$ is unramified over $B$ outside $S$, so $e_K(\mathfrak p)=1$,
and a prime $\mathfrak P$ of $K$ above $\mathfrak p$ has a Frobenius
$\Frob_{\mathfrak P}\in\Gal(K/B)$ of order $f_K(\mathfrak p)$. Choose
$g\in G_B$ mapping to $\Frob_{\mathfrak P}$. Its image in
$G_B/D_2G_B=\Gal(E/B)$ is the Frobenius of $\mathfrak p$ in $E$
(Lemma~\ref{gf:in-tower}), whose elementary image is $\overline{\Frob}_{\mathfrak p}$; it is
nonzero, since $0\in\Sigma_2$. By Lemma~\ref{tw:quadratic-layer}, which uses
the complete list of generators of $N_B$ in
Lemma~\ref{tw:GB-presentation}, and since $\overline{\Frob}_{\mathfrak p}^{[2]}\notin R_2$,
$g$ has order at least $4$ in $G_B/D_3G_B$. The projection $G_B\to G_B/D_3G_B$
factors through $\Gal(K/B)$, because $K$ contains the fixed field of
$D_4G_B\subseteq D_3G_B$. So $\Frob_{\mathfrak P}$ has order at least $4$, and
this order is a power of two because $\Gal(K/B)$ is a $2$-group.
\end{proof}


\begin{lemma}\label{an:sigma2}
The set $\Sigma_2$ of Definition~\ref{an:P4} has exactly $21$ elements: zero
and the following twenty vectors, written as strings $c_0c_1\cdots c_7$ of
their coordinates, as in Section~\ref{tw:tower}:
\[
\begin{array}{ccccc}
 00100000&11110010&11010010&00010000&10000001\\
 10010001&00001000&10100111&10101111&00000100\\
 11101011&11101111&00000010&01100101&01100111\\
 00000001&01011010&01011011&10111010&11000101.
\end{array}
\]
In the notation of Table~\ref{tw:vector-table}, these twenty vectors are the
elementary images of the local involutions $x_j$, $y_j$ and $x_jy_j$ for
$j=1,2$, of the nonidentity elements of the four local groups
$C_2\times C_2$ above $3$ and $5$, and of $\iota_1$ and $\iota_2$.
\end{lemma}

\begin{proof}
Each of these images lies in $\Sigma_2$ by Lemma~\ref{tw:quadratic-layer},
since an involution has trivial square. The supplementary program
\texttt{census241.py} computes $\Sigma_2$ by linear algebra over $\Ftwo$
from the relations of Lemma~\ref{tw:quadratic-layer}, and finds no other
vectors.
\end{proof}

\begin{proposition}[Frobenius vectors of the primes of norm at most $10^6$]\label{an:census-prop}
Among the $78616$ free primes
$\mathfrak p$ of $B$ with $\mathrm N\mathfrak p\le10^6$, exactly $246$ have
$\overline{\Frob}_{\mathfrak p}=0$, exactly $5994$ have
$\overline{\Frob}_{\mathfrak p}\in\Sigma_2\setminus\{0\}$, and the remaining $72376$ form
$\mathcal P_4$.
\end{proposition}

\begin{proof}
The supplementary program \texttt{census241.py} (directory
\path{papers/0.04273/certificates} of \cite{NaslundCode4317}) computes the
Frobenius vector of each of these primes from
Proposition~\ref{gf:local-data}(1) with exact modular arithmetic, and compares
it with the elements of $\Sigma_2$ (Lemma~\ref{an:sigma2}). It was written for
the variant of the construction with the cap at $\mathfrak t_0$ instead of
$\mathfrak t_3$ (Remark~\ref{tw:same-quotient}), whose free primes differ
from the present ones only in that $\mathfrak t_0$ is free and $\mathfrak t_3$
is not; both have Frobenius vectors outside $\Sigma_2$, so the three counts
are the same for both. The program \texttt{adaptive41.py} enumerates the free
primes of the present tower and reproduces the three counts.
\end{proof}

Below we use this classification of the primes, not the three counts.


\begin{remark}\label{an:base-ceiling}
With $E$ as the comparison field, Proposition~\ref{an:prime-budget} at
$\sigma=301/300$ gives the bound $C=\ceiling$ for the present tower. Its
lower bounds
$(e^0_{\mathfrak p},f^0_{\mathfrak p})$ are the types in $K$ at the selected
primes, $(1,4)$ at the primes of $\mathcal P_4$
(Lemma~\ref{an:census-degree}), and the types in $E$ at the other free
primes. By \eqref{an:local-sum} and Lemma~\ref{an:local-contribution} one may
take $Y_*=Y_E(\sigma)-\Delta^E_{\rm sel}(\sigma)-\Delta^E_4(\sigma)$, where
$\Delta^E_{\rm sel}(\sigma)$ is the sum over the selected primes of the
difference between the terms in \eqref{an:local-sum} for $M=E$ and for
$M=K$, and
$\Delta^E_4(\sigma)=\sum_{\mathfrak p\in\mathcal P_4}\bigl(g_\sigma(\mathrm
N\mathfrak p^2)/4-g_\sigma(\mathrm N\mathfrak p^4)/8\bigr)$; $Y'_*$ is bounded
from the same lower bounds. Compared with the variant with the cap at
$\mathfrak t_0$ (Remark~\ref{tw:same-quotient}), the terms
of $\mathfrak t_0$ and of $\mathfrak t_3$ change places between $\Delta^E_{\rm sel}$ and
$\Delta^E_4$, with the same formula
$g_\sigma(\mathrm N\mathfrak p^2)/4-g_\sigma(\mathrm N\mathfrak p^4)/8$, so the
totals are the same for both; the supplementary program \texttt{ceiling41.py}
(directory \path{papers/0.042901/certificates} of \cite{NaslundCode4317})
computes the upper bound $0.04871284289\ldots$. With this
$C$, the margin of Section~\ref{sec:certificate} with re-optimized profiles
gives the exponent $1.042901$ (Table~\ref{intro:table-steps}). This bound is
not used in the proof of Theorem~\fref{thm:main}.
\end{remark}

\begin{remark}\label{an:limit-remark}
The selected primes alone limit what Proposition~\ref{an:signed} can give.
In the notation of its proof, $\Sigma_{\mathsf K}\ge0$ by condition (ii), and
$2\kappa_\infty-T_{\rm sel}>1.07>0$. So no system of floors, abscissae and
coefficients gives a bound below
$Z_{\rm sel}(1)=0.04166848403\ldots$ (Lemma~\ref{an:selected-lemma}), the
contribution of the selected primes to $Z(1)$. It exceeds the corresponding
number for the variant with the cap at $\mathfrak t_0$ by $2.3\cdot10^{-8}$,
because the selected primes above $41$ replace those above $7$. The bound $C_{\rm eff}=\Ceff$
exceeds $Z_{\rm sel}(1)$ by $6.08\cdot10^{-4}$; Section~\ref{lim:section}
discusses what this means for the exponent.
\end{remark}

\begin{remark}[The Lean formalization]\label{an:lean-hypothesis}
The Lean formalization \cite{NaslundLeanW} proves a weaker form of
Theorem~\fref{thm:main}, under one explicit hypothesis that it does not
prove: for the fields of the present tower, there are finite planar sets
$U_j$ with $|U_j|\to\infty$ and $u(U_j)/|U_j|^{20863/20000}\to\infty$, where
$20863/20000=1.04315$, provided that
\begin{equation}\label{an:hw}
 \frac{\log\zeta_{\EW}(1+\frac1{4411})}{8192}
 +\frac1{4411}\Bigl(\frac{\ell-\gamma-\log(4\pi)}4
 -\frac1{8192}\frac{\zeta_{\EW}'}{\zeta_{\EW}}(2)\Bigr)<0.050969,
\end{equation}
where $\EW$ is the field of Definition~\ref{dh:W}, given in the
formalization by explicit generators $\sqrt{\beta_i}$ over $E$ for four
dihedral forms that span $W''$. This is the bound of
Proposition~\ref{an:prime-budget} at $\sigma=1+1/4411$ with comparison field
$\EW$ in place of $E$. The formalization proves in Lean that
$[\EW:\Q]=8192$, that $\EW$ lies in the fixed base of the tower, the local
data of $\EW$ at $2$ and at fifty split primes below $1000$, the resulting
bound $0.0425$ for $d^{-1}\log L_F(1)$, and the margin at the exponent
$1.04315$. The supplementary program \texttt{h\_w\_receipt.py} bounds the
left side of \eqref{an:hw} by $0.0509172$, from the enclosures of
Proposition~\fref{lv:values} and an explicit bound for the logarithmic
derivative. The formalization does not cover the census beyond norm $1000$,
the signed kernel or the general shell weights.
\end{remark}

